\subsubsection{Editorial-Oriented Two-Stage Paradigm}

A linear TF--IDF baseline using word- and character-level n-grams, source
indicators, and simple numeric features is first evaluated. Despite extensive
hyperparameter exploration (Table~\ref{tab:baseline_hyperparams}), performance
saturates around a macro-averaged F1 score of $0.719$, indicating that a
single-stage linear formulation largely exhausts the available signal.

\begin{table}[H]
\centering
\caption{Optimal hyperparameter ranges for the linear TF--IDF baseline.}
\label{tab:baseline_hyperparams}
\begin{tabular}{lc}
\hline
\textbf{Hyperparameter} & \textbf{Selected value / range} \\
\hline
Word n-grams        & 2 \\
Character n-grams   & 5--7 \\
Regularization $C$  & 0.6--1.0 \\
Max document freq.  & 0.85--0.9 \\
\hline
\end{tabular}
\end{table}

Inspection of feature weights reveals that the strongest predictors are
editorial and structural cues, such as publisher identifiers, RSS markers,
and HTML fragments. Several of these tokens exhibit near-deterministic class
associations, as illustrated in Table~\ref{tab:low_entropy_tokens}.

\begin{table}[H]
\centering
\caption{Examples of high-support editorial tokens with near-zero conditional entropy.}
\label{tab:low_entropy_tokens}
\small
\resizebox{\linewidth}{!}{%
\begin{tabular}{lcccl}
\hline
\textbf{Token} & \textbf{Type} & \textbf{Support} & \textbf{Entropy} & \textbf{Dominant label} \\
\hline
\texttt{<h4>} & editorial & 628 & 0.00 & 2 \\
\texttt{</h4>} & editorial & 628 & 0.00 & 2 \\
\texttt{<p align="right">} & editorial & 311 & 0.00 & 2 \\
\texttt{<a href=...>} & editorial & 233 & 0.00 & 2 \\
\texttt{<img class="rss-ad">} & editorial & 196 & 0.00 & 2 \\
\texttt{<strong>} & editorial & 111 & 0.00 & 2 \\
\texttt{<cite>} & editorial & 40 & 0.00 & 2 \\
\texttt{<em>} & editorial & 57 & 0.47 & 2 \\
\hline
\end{tabular}}
\end{table}

To formalize this behavior, for a token $t$ and class label $Y$ we estimate the
empirical distribution $p(c \mid t)$ on the development set as the fraction of
documents containing $t$ that belong to class $c$. The conditional entropy is
defined as:
\[
H(Y \mid t) = - \sum_{c \in \mathcal{Y}} p(c \mid t)\log p(c \mid t).
\]
For a non-trivial subset of tokens, $H(Y \mid t)=0$, meaning that their presence
alone deterministically identifies the class in the development set.

Based on these observations, we adopt a Two-Stage classification paradigm. For
each token $t$, we compute:
\[
\text{purity}(t) = \max_{c \in \mathcal{Y}} p(c \mid t), \quad
\text{support}(t) = |\{x_i : t \in x_i\}|,
\]
where $\text{support}(t)$ denotes the number of development documents
containing $t$. Tokens satisfying purity and support thresholds $\tau_p$ and
$\tau_s$ are collected into a rule set:
\[
\mathcal{R} = \{t \in \mathcal{T} : \text{purity}(t) \ge \tau_p \;\wedge\;
\text{support}(t) \ge \tau_s\}.
\]

\paragraph{Stage 1: Deterministic Assignment}
If a document $x$ contains any token $t \in \mathcal{R}$, it is assigned the
corresponding dominant class:
\[
\hat{y}^{(1)}(x) = \arg\max_{c \in \mathcal{Y}} p(c \mid t).
\]

\paragraph{Stage 2: Probabilistic Classification}
All remaining documents are classified by a probabilistic model $f_\theta(x)$:
\[
\hat{y}^{(2)}(x) = \arg\max_{c \in \mathcal{Y}} f_\theta(x)_c.
\]

\paragraph{Unified Decision Rule}
\[
\hat{y}(x) =
\begin{cases}
\hat{y}^{(1)}(x), & \text{if } \exists\, t \in \mathcal{R} \text{ such that } t \in x, \\
\hat{y}^{(2)}(x), & \text{otherwise}.
\end{cases}
\]

This formulation explicitly separates deterministic editorial patterns from
semantically ambiguous instances, while reducing to a standard single-stage
classifier when no high-purity rules apply.
